\documentclass[10pt,a4paper]{article}
% Build from the repository root with:
% tectonic -X compile docs/ilc_dpd_forward_model_theory.tex --outdir output/pdf

\usepackage[margin=22mm,headheight=20pt,headsep=7mm,footskip=11mm]{geometry}
\usepackage{microtype}
\usepackage{fontspec}
\usepackage{amsmath}
\usepackage{unicode-math}
\usepackage{enumitem}
\usepackage[dvipsnames]{xcolor}
\usepackage{booktabs}
\usepackage{fancyhdr}
\usepackage{titlesec}
\usepackage{xurl}
\usepackage{hyperref}

\setmainfont{texgyretermes-regular.otf}[
  BoldFont=texgyretermes-bold.otf,
  ItalicFont=texgyretermes-italic.otf,
  BoldItalicFont=texgyretermes-bolditalic.otf
]
\setmathfont{texgyretermes-math.otf}

\definecolor{TitleBlue}{HTML}{17365D}
\definecolor{LinkBlue}{HTML}{1F5A94}
\definecolor{RuleGray}{HTML}{B8C2CC}

\hypersetup{
  colorlinks=true,
  linkcolor=LinkBlue,
  urlcolor=LinkBlue,
  citecolor=LinkBlue,
  pdftitle={From Scalar Residual Injection to Iteration-Wise PA-Model Adjoint Updates},
  pdfauthor={Remote DPD Project},
  pdfsubject={A mathematical note on waveform ILC for digital predistortion},
  pdfkeywords={digital predistortion, iterative learning control, PA model, VJP, Gauss-Newton, Levenberg-Marquardt}
}

\pagestyle{fancy}
\fancyhf{}
\fancyhead[L]{\small\color{TitleBlue}Forward-Model Backpropagation for ILC DPD}
\fancyhead[R]{\small\color{TitleBlue}Mathematical Note}
\fancyfoot[C]{\small\thepage}
\renewcommand{\headrulewidth}{0.4pt}
\renewcommand{\headrule}{\hbox to\headwidth{\color{RuleGray}\leaders\hrule height \headrulewidth\hfill}}

\titleformat{\section}
  {\large\bfseries\color{TitleBlue}}
  {\thesection}{0.65em}{}
\titleformat{\subsection}
  {\normalsize\bfseries\color{TitleBlue}}
  {\thesubsection}{0.65em}{}
\titlespacing*{\section}{0pt}{1.6em}{0.55em}
\titlespacing*{\subsection}{0pt}{1.2em}{0.4em}

\setlength{\parindent}{0pt}
\setlength{\parskip}{0.55em plus 0.1em minus 0.1em}
\setlist[itemize]{leftmargin=1.6em,itemsep=0.25em,topsep=0.35em}
\setlist[enumerate]{leftmargin=1.8em,itemsep=0.3em,topsep=0.35em}
\allowdisplaybreaks[2]
\raggedbottom
\clubpenalty=10000
\widowpenalty=10000

\newcommand{\Rinner}[2]{\left\langle #1,#2\right\rangle_{\mathbb R}}
\newcommand{\norm}[1]{\left\lVert #1\right\rVert_2}
\newcommand{\rms}{\operatorname{RMS}}
\newcommand{\roll}{\operatorname{roll}}
\newcommand{\Fhat}{\widehat F}
\newcommand{\Jhat}{\widehat J}

\title{\vspace{-1.2cm}\color{TitleBlue}\bfseries
  From Scalar Residual Injection to Iteration-Wise\\[0.18em]
  PA-Model Adjoint Updates}
\author{\large A Mathematical Note on Waveform Digital Predistortion}
\date{August 2026}

\begin{document}

\maketitle
\vspace{-1.2em}
\begin{center}
  \color{RuleGray}\rule{0.84\textwidth}{0.7pt}
\end{center}

\begin{abstract}
This note formulates waveform iterative learning control (ILC) for digital predistortion (DPD) as a
constrained inverse problem. It shows that scalar residual injection implicitly replaces the local
power-amplifier (PA) sensitivity by a scaled identity, derives the resulting convergence and descent
limitations, and then formulates an iteration-wise forward-model alternative. The measured residual is
backpropagated through a frozen real-linear PA-model adjoint; a damped Gauss-Newton or
Levenberg-Marquardt step supplies regularized local inversion. The presentation is theoretical and
contains no experimental claims.
\end{abstract}

\tableofcontents
\vspace{0.4em}

\section{Scope and notation}

Throughout this note, \emph{classical scalar ILC} denotes the waveform update that directly injects a
scaled output residual into the PA input. This terminology does not cover the full lifted-domain ILC
literature, in which a designed learning operator, robustness filter, identified plant inverse, or
instantaneous-gain preconditioner may be used. The limitation studied here is therefore not repetition-domain
learning itself, but the scalar identity-sensitivity approximation
\cite{chani2016ilc,schoukens2017preinverse}.

Let
\begin{itemize}
  \item $u\in\mathbb C^N$ be the finite waveform applied to the PA;
  \item $F:\mathbb C^N\rightarrow\mathbb C^N$ be the calibrated PA input-output operator;
  \item $d\in\mathbb C^N$ be the desired PA output;
  \item $y_k=F(u_k)$ be the measured output at outer iteration $k$; and
  \item $r_k=d-y_k$ be the output tracking residual.
\end{itemize}

Some implementation descriptions instead use $e_k=y_k-d=-r_k$. This note retains $r_k=d-y_k$
throughout, so the classical update has the familiar plus-residual form.

The waveform-level inverse problem is
\begin{equation}
  \min_{u\in\mathcal U}\;\mathcal L(u),
  \qquad
  \mathcal L(u)
  =\frac12\norm{F(u)-d}^{2}
  =\frac12\norm{r(u)}^{2},
  \label{eq:inverse-problem}
\end{equation}
where $\mathcal U$ is the feasible set imposed by peak, RMS, PAPR, and hardware limits.

The derivation assumes a fixed input/output coordinate system. If the measured output is independently
renormalized at every iteration, the effective operator is an iteration-dependent composition of the PA and
the calibration map rather than the physical PA map alone.

A nonlinear RF PA is generally not holomorphic as a complex map because its basis functions depend on
$|u|$. We therefore treat $\mathbb C^N$ as $\mathbb R^{2N}$, with real inner product
\begin{equation}
  \Rinner{a}{b}=\operatorname{Re}(a^H b).
  \label{eq:real-inner-product}
\end{equation}
Let $J(u)$ denote the real-linear Fr\'echet derivative of $F$:
\begin{equation}
  F(u+h)=F(u)+J(u)h+o(\norm{h}).
\end{equation}
Its real adjoint $J(u)^T$ is defined by
\begin{equation}
  \Rinner{v}{J(u)h}=\Rinner{J(u)^T v}{h}.
  \label{eq:real-adjoint}
\end{equation}
Because $r(u)=d-F(u)$, the directional derivative of the loss is
\begin{equation}
  D\mathcal L(u)[h]
  =-\Rinner{r(u)}{J(u)h}
  =\Rinner{-J(u)^T r(u)}{h}.
\end{equation}
Hence the true input-space gradient is
\begin{equation}
  \boxed{\nabla\mathcal L(u)=-J(u)^T r(u).}
  \label{eq:true-gradient}
\end{equation}
The output residual is therefore not, in general, an input correction. It must be transported through the
PA adjoint before it becomes a gradient in the input coordinates.

\section{Classical scalar ILC as an identity-sensitivity approximation}

The scalar linear ILC update is
\begin{equation}
  u_{k+1}=u_k+\mu r_k,
  \qquad \mu>0.
  \label{eq:classical-ilc}
\end{equation}
Using a first-order PA expansion at $u_k$ gives
\begin{align}
  r_{k+1}
  &=d-F(u_k+\mu r_k) \notag\\
  &=r_k-\mu J_k r_k
    +o\!\left(\mu\norm{r_k}\right),
  \label{eq:local-residual-dynamics}
\end{align}
where $J_k=J(u_k)$. Thus
\begin{equation}
  r_{k+1}\approx (I-\mu J_k)r_k.
\end{equation}
Equation~\eqref{eq:classical-ilc} implicitly treats the output residual as if it were already expressed in
the PA input coordinates. Equivalently, it approximates the local inverse learning operator, or the
adjoint sensitivity after step scaling, by a scalar multiple of the identity.

\subsection{Linear and scalar convergence conditions}

For a fixed linear PA $F(u)=Au$, asymptotic convergence requires
\begin{equation}
  \rho(I-\mu A)<1,
  \label{eq:spectral-condition}
\end{equation}
where $\rho(\cdot)$ is the spectral radius. For a nonzero scalar complex gain $A=g$, a real positive learning rate
must satisfy
\begin{equation}
  |1-\mu g|<1
  \quad\Longleftrightarrow\quad
  0<\mu<\frac{2\operatorname{Re}(g)}{|g|^2}.
  \label{eq:complex-gain-condition}
\end{equation}
Such a $\mu$ exists only when $\operatorname{Re}(g)>0$. A sufficiently large PA phase rotation can therefore
turn the nominal correction into a non-contractive direction. For example, if $g=\mathrm j$, then
$|1-\mathrm j\mu|>1$ for every $\mu>0$.

The same conclusion follows directly from the loss derivative. Along the classical step $h=\mu r_k$,
\begin{equation}
  D\mathcal L(u_k)[\mu r_k]
  =-\mu\Rinner{r_k}{J_k r_k}.
  \label{eq:classical-descent-test}
\end{equation}
The step is a first-order descent direction only if
\begin{equation}
  \Rinner{r_k}{J_k r_k}>0.
  \label{eq:descent-condition}
\end{equation}
This condition is not guaranteed for a general complex, nonlinear, or dynamic PA.

\subsection{Main limitations}

\begin{enumerate}
  \item \textbf{Direction error.} Phase distortion or negative local AM/AM slope can make
  $\Rinner{r_k}{J_k r_k}\leq 0$. Identity transport can then be orthogonal to, or point against, the true
  descent direction.

  \item \textbf{Scale error.} If the local PA slope is small, the output changes little although $\mu r_k$
  continues to accumulate at the input. If the slope is large, the same fixed $\mu$ can overshoot. One scalar
  cannot simultaneously invert widely different singular directions.

  \item \textbf{Memory and cross-sample coupling.} With PA memory, $J_k$ is not diagonal. An error at one
  time index may require corrections at several input indices. Identity transport ignores this coupling.

  \item \textbf{Iteration-dependent nonlinearity.} The Jacobian changes with $u_k$. A learning rate that is
  stable in one envelope region need not remain stable after the waveform moves to another region.

  \item \textbf{Saturation and infeasibility.} In deep saturation, $J_k$ may become rank deficient or develop
  very small singular values. Scalar ILC still injects $\mu r_k$ along locally insensitive directions. If
  $d\notin F(\mathcal U)$, then, for compact $\mathcal U$ and continuous $F$,
  \begin{equation}
    \min_{u\in\mathcal U}\norm{F(u)-d}>0.
    \label{eq:unreachable-target}
  \end{equation}
  No ILC law can drive such an error to zero.

  \item \textbf{Noise integration.} If $y_k=F(u_k)+n_k$, the classical update directly contains
  $-\mu n_k$. Repetition can accumulate capture noise unless averaging, filtering, regularization, or a
  stopping rule is used.
\end{enumerate}

These are limitations of the identity-Jacobian baseline, not impossibility statements about every form of
ILC. A known plant inverse, instantaneous complex gain, or frequency-dependent learning filter can correct
some of them when its structural assumptions hold.

\subsection{Instantaneous gain as a diagonal secant preconditioner}

For nonzero input samples, instantaneous-gain ILC estimates a pointwise secant gain
\begin{equation}
  \widehat g_k[n]=\frac{y_k[n]}{u_k[n]},
  \qquad
  \widehat D_k=\operatorname{diag}(\widehat g_k[n]).
  \label{eq:instantaneous-gain}
\end{equation}
A damped diagonal update is
\begin{equation}
  \delta_k
  =\left(\widehat D_k^H\widehat D_k+\lambda I\right)^{-1}
   \widehat D_k^H r_k,
  \qquad \lambda>0.
  \label{eq:instantaneous-update}
\end{equation}
This is a diagonal secant preconditioner, not the full differential inverse of a general nonlinear PA. To
see the distinction, let $f(u)=a(\varrho)u$ with $\varrho=|u|>0$. Its real-linear differential is
\begin{equation}
  Df(u)[h]
  =\left(a(\varrho)+\frac{\varrho}{2}a'(\varrho)\right)h
  +\frac{a'(\varrho)u^2}{2\varrho}\,\overline h.
  \label{eq:secant-versus-differential}
\end{equation}
The ratio $y/u=a(\varrho)$ omits the amplitude-dependent differential terms, the conjugate path, and all
cross-sample memory coupling. Low-amplitude division also requires thresholding. Nevertheless,
instantaneous gain can be highly effective for approximately memoryless multiplicative distortion and is not
assumed to be weaker than the forward-model method \cite{chani2016ilc,wei2026icg}.

\section{Iteration-wise forward PA modeling}

Backpropagation through a fitted forward PA model has established precedents
\cite{tarver2019backprop,loebl2023memorybackprop}. In the waveform-level construction considered here, the
model-adjoint scheme estimates the missing local PA geometry from the latest measured pair $(u_k,y_k)$ and
refits a differentiable forward model at every outer ILC iteration:
\begin{equation}
  \Fhat_k(u)=\Fhat(u;\widehat\theta_k),
\end{equation}
for example by regularized regression,
\begin{equation}
  \widehat\theta_k
  =\arg\min_\theta
  \left\|\Fhat(u_k;\theta)-y_k\right\|_2^2
  +\gamma\norm{\theta}^{2}.
  \label{eq:model-fit}
\end{equation}
The model regularization $\gamma$ is distinct from the update damping introduced later.

The fitted parameters $\widehat\theta_k$ and envelope scale are frozen while iteration $k$'s update is
computed. Gradients are not propagated through the regression solver, measurement process, envelope scale,
delay estimator, or calibration estimator. The model is refitted only after a new PA input-output pair has
been observed.

This yields two nested loops:
\begin{enumerate}
  \item an \emph{identification loop}, which estimates a local forward PA model from the current measurement;
  \item a \emph{learning loop}, which uses the frozen model derivative to transport the measured residual
  back to the PA input.
\end{enumerate}
The residual remains the physical measurement residual $r_k=d-y_k$; it is not replaced by a model residual.

\section{Real-linear backpropagation through the PA model}

Let
\begin{equation}
  \Jhat_k=D\Fhat_k(u_k)
\end{equation}
be the frozen model derivative. Its Jacobian-vector product (JVP) maps an input perturbation to an output
perturbation, $h\mapsto\Jhat_kh$, and its vector-Jacobian product (VJP) transports an output cotangent back
to the input, $v\mapsto\Jhat_k^Tv$.

\subsection{Memory-polynomial example}

Consider a circular memory polynomial \cite{morgan2006gmp}
\begin{equation}
  \Fhat_k(u)[n]
  =\sum_{m=0}^{M-1}\sum_{p\in\mathcal P}
  c_{p,m}\,z_m[n]
  \left(\frac{|z_m[n]|}{s}\right)^{p-1},
  \qquad
  z_m[n]=u[(n-m)\bmod N].
  \label{eq:memory-polynomial}
\end{equation}
Here $\mathcal P\subset\{1,3,5,\ldots\}$ is a set of positive odd orders. For one basis function
\begin{equation}
  \phi_p(z)=c\,z\left(\frac{|z|}{s}\right)^{p-1},
\end{equation}
the real-linear differential contains both direct and conjugate terms:
\begin{equation}
  D\phi_p(z)[h]=a_p(z)h+b_p(z)\overline h,
  \label{eq:basis-differential}
\end{equation}
where
\begin{align}
  a_p(z)
  &=c\frac{p+1}{2}
    \left(\frac{|z|}{s}\right)^{p-1}, \\
  b_p(z)
  &=c\frac{p-1}{2}
    \left(\frac{z}{s}\right)^2
    \left(\frac{|z|}{s}\right)^{p-3},
  \qquad p>1.
\end{align}
For $p=1$, define $b_1(z)=0$ explicitly; at $z=0$, use the continuous extension $b_p(0)=0$. After summing
polynomial orders at each delay into arrays $a_m$ and $b_m$, the JVP is
\begin{equation}
  \Jhat_k h
  =\sum_m\left[
    a_m\odot\roll(h,m)
    +b_m\odot\overline{\roll(h,m)}
  \right].
  \label{eq:jvp}
\end{equation}
The corresponding real adjoint is
\begin{equation}
  \Jhat_k^T v
  =\sum_m\roll\!\left(
    \overline{a_m}\odot v+b_m\odot\overline v,-m
  \right).
  \label{eq:vjp}
\end{equation}
The conjugate contribution in Eq.~\eqref{eq:vjp} is $b_m\odot\overline v$, not
$\overline{b_m}\odot\overline v$. An ordinary complex-linear Hermitian derivative would omit this path and generally
produce the wrong gradient.

\section{Backpropagated learning laws}

\subsection{Raw model-adjoint update}

Replacing identity transport by the learned adjoint gives
\begin{equation}
  \delta_k^{\mathrm{VJP}}=\eta\Jhat_k^T r_k,
  \qquad
  u_{k+1}=u_k+\delta_k^{\mathrm{VJP}},
  \qquad \eta>0.
  \label{eq:raw-vjp}
\end{equation}
If $\Jhat_k=I$, this reduces exactly to scalar linear ILC with $\eta=\mu$. With an exact derivative,
$\Jhat_k=J_k$, the first-order loss change is
\begin{equation}
  D\mathcal L(u_k)[\delta_k^{\mathrm{VJP}}]
  =-\eta\norm{J_k^T r_k}^{2}<0
  \label{eq:exact-vjp-descent}
\end{equation}
whenever the true gradient is nonzero. This negative directional derivative guarantees actual decrease only
for a sufficiently small finite step; it is not a global or arbitrary-step guarantee.

With derivative error $J_k=\Jhat_k+E_k$,
\begin{equation}
  D\mathcal L(u_k)[\delta_k^{\mathrm{VJP}}]
  =-\eta\Rinner{J_k^T r_k}{\Jhat_k^T r_k}.
\end{equation}
Here $\|E_k\|_2$ is the operator 2-norm induced by the $2N\times2N$ real representation of the real-linear
map $E_k$.
Furthermore,
\begin{align}
  \Rinner{J_k^T r_k}{\Jhat_k^T r_k}
  &\geq \norm{\Jhat_k^T r_k}
  \left(
    \norm{\Jhat_k^T r_k}-\|E_k\|_2\norm{r_k}
  \right).
  \label{eq:model-error-bound}
\end{align}
A sufficient local descent condition is therefore
\begin{equation}
  \|E_k\|_2\norm{r_k}<\norm{\Jhat_k^T r_k}.
  \label{eq:model-quality-condition}
\end{equation}
The learned adjoint is thus a surrogate gradient. Backpropagation through a fitted model does not by itself
establish that the direction equals the true PA gradient.

\subsection{Damped Gauss-Newton / Levenberg-Marquardt update}

Raw gradient descent still inherits poor Jacobian scaling. A regularized local step instead solves
\begin{equation}
  \delta_k
  =\arg\min_\delta
  \frac12\norm{r_k-\Jhat_k\delta}^{2}
  +\frac{\lambda_k}{2}\norm{\delta}^{2},
  \qquad \lambda_k>0.
  \label{eq:lm-objective}
\end{equation}
Its first-order condition is the damped normal equation
\begin{equation}
  \boxed{
  \left(\Jhat_k^T\Jhat_k+\lambda_k I\right)\delta_k
  =\Jhat_k^T r_k.}
  \label{eq:lm-normal-equation}
\end{equation}
The normal operator is positive definite in the real geometry because, for every nonzero $q$,
\begin{equation}
  \Rinner{q}{(\Jhat_k^T\Jhat_k+\lambda_k I)q}
  =\norm{\Jhat_kq}^{2}+\lambda_k\norm{q}^{2}>0.
  \label{eq:normal-spd}
\end{equation}
Consequently, the step can be computed by matrix-free conjugate gradients using only
\begin{equation}
  q\longmapsto\Jhat_kq
  \longmapsto\Jhat_k^T(\Jhat_kq)+\lambda_kq,
\end{equation}
without constructing the full $2N\times2N$ real Jacobian. The same frozen linearization must be used within
one CG solve so that the normal operator remains symmetric positive definite.

If $\Jhat_k=U\Sigma V^T$ is the singular-value decomposition of its real representation, then
\begin{equation}
  \delta_k
  =V\,\operatorname{diag}\!\left(
    \frac{\sigma_i}{\sigma_i^2+\lambda_k}
  \right)U^T r_k.
  \label{eq:lm-svd}
\end{equation}
Well-observed directions are approximately inverted, whereas weak directions are suppressed. Indeed,
\begin{equation}
  0\leq\frac{\sigma}{\sigma^2+\lambda_k}
  \leq\frac{1}{2\sqrt{\lambda_k}}.
\end{equation}

Two limiting cases clarify the connection to simpler ILC laws:
\begin{itemize}
  \item If $\Jhat_k=I$, then $\delta_k=r_k/(1+\lambda_k)$, a damped scalar ILC step.
  \item For a memoryless complex gain $g$,
  \begin{equation}
    \delta_k=\frac{g^*}{|g|^2+\lambda_k}r_k,
  \end{equation}
  which is a Tikhonov-regularized complex-gain inverse.
\end{itemize}

For a perfect linear model $F(u)=Au$, the new residual is
\begin{equation}
  r_{k+1}
  =\left[I-A(A^TA+\lambda I)^{-1}A^T\right]r_k.
\end{equation}
Along each nonzero singular direction, the contraction factor is
\begin{equation}
  \frac{\lambda}{\sigma_i^2+\lambda}.
\end{equation}
With $\lambda=0$ and invertible $A$, the exact linear inverse is recovered in one idealized step.

\section{Measurement-anchored prediction and safeguards}

A forward model can have static bias at the current waveform. Using $\Fhat_k(u_k+\delta)$ directly would
mistake that bias for a measured output change. Instead use the anchored prediction
\begin{equation}
  \widetilde y_k(\delta)
  =y_k+\Fhat_k(u_k+\delta)-\Fhat_k(u_k).
  \label{eq:anchored-prediction}
\end{equation}
It satisfies $\widetilde y_k(0)=y_k$ exactly while retaining the model-predicted nonlinear increment. A
candidate can be accepted only when, for example,
\begin{equation}
  \frac12\norm{d-\widetilde y_k(\delta)}^{2}
  <\frac12\norm{r_k}^{2}.
  \label{eq:anchored-acceptance}
\end{equation}
This is a predicted-decrease condition under the model, not a guarantee of monotone decrease for the
physical PA.

The update should additionally obey a trust region and the feasible input set:
\begin{equation}
  \rms(\delta)\leq\tau\rms(u_k),
  \qquad
  u_k+\delta\in\mathcal U.
  \label{eq:safeguards}
\end{equation}
If the full step is rejected, test $\beta^b\delta$, $0<\beta<1$, until a candidate passes the anchored
decrease and input constraints. If none passes, hold $u_k$. A model-fit failure should likewise produce an
explicit hold or declared scalar-ILC fallback rather than silently reusing an unrelated model.

These safeguards bound the consequences of local model error. They do not create an inverse in saturation,
guarantee global convergence, or ensure branch switching in a nonmonotone PA.

\section{Saturation, reachability, and remaining limitations}

The learned model changes the update from identity transport,
\begin{equation}
  \delta_k=\mu r_k,
\end{equation}
to model-adjoint transport,
\begin{equation}
  \delta_k=\eta\Jhat_k^Tr_k,
\end{equation}
or a regularized local inverse, Eq.~\eqref{eq:lm-normal-equation}. This directly addresses phase rotation,
unequal local slope, and memory coupling when the forward model is locally accurate. Iterative refitting lets
the fitted derivative track the waveform through the PA operating region.

The method remains local and model-dependent:
\begin{itemize}
  \item it has no global convergence guarantee for a nonconvex PA inverse;
  \item model mismatch can rotate the learned gradient away from the true gradient;
  \item a poorly excited waveform cannot identify unobserved model directions;
  \item changing calibration across iterations changes the effective operator being learned;
  \item capture noise affects both the residual and the fitted model;
  \item the optimized variable is one finite, repeatable waveform, so arbitrary traffic requires a separate
  parameterized DPD identification stage; and
  \item if $d\notin F(\mathcal U)$, no gradient or second-order approximation can make the target reachable.
\end{itemize}

In saturation, $J_k$ may become rank deficient. A residual component in $\ker(J_k^T)$ cannot be reduced to
first order. In the extreme locally constant idealization $J_k=0$, if the learned model faithfully represents
that behavior, then $\Jhat_k=0$ and the right-hand side of Eq.~\eqref{eq:lm-normal-equation} vanishes. Positive
damping gives $\delta_k=0$. This is a safe indication of local unresponsiveness, not recovery of the
unreachable desired waveform. Conversely, a small learned gradient can also result from model mismatch, so a
local saturation indicator is not a proof of global unreachability.

\section{Compact algorithm and interpretation}

At each outer iteration $k$:
\begin{enumerate}
  \item apply $u_k$, measure and align $y_k$, and form $r_k=d-y_k$;
  \item fit and validate $\Fhat_k$ from the current input-output pair;
  \item freeze $\Fhat_k$ and construct its real-linear JVP and VJP at $u_k$;
  \item compute a raw VJP step or solve Eq.~\eqref{eq:lm-normal-equation} by matrix-free CG;
  \item apply trust-region scaling, input projection, anchored nonlinear prediction, and backtracking;
  \item accept the first safe predicted-decrease candidate, otherwise hold or use a declared fallback; and
  \item refit only after the next physical PA measurement.
\end{enumerate}

The conceptual improvement is not merely the addition of a model. It replaces the identity-sensitivity
approximation with the differential of an iteration-wise fitted forward model, uses the real adjoint to
transport the measured residual to the PA input, and regularizes that transport before the next ILC experiment.

This construction combines established PA modeling, backpropagation, and ILC ingredients in a specific
waveform-level loop. It does not claim the first use of PA forward models, PA Jacobians, memory
backpropagation, direct learning, memory-polynomial modeling, or model-based ILC. Here ``improved'' means
that the update contains richer local sensitivity information and explicit numerical safeguards; it does not
mean uniform superiority over instantaneous-gain, BLA-inverse, or other plant-informed ILC methods.

\section*{Companion documents}
\addcontentsline{toc}{section}{Companion documents}

The following repository documents provide the implementation and prior-art context:
\begin{itemize}
  \item algorithm conventions: \path{docs/algorithm_design.md};
  \item prior-art boundaries: \path{docs/prior_art_review.md}; and
  \item Markdown edition: \path{docs/ilc_dpd_forward_model_theory.md}.
\end{itemize}

\addcontentsline{toc}{section}{References}
\begin{thebibliography}{9}
\small
\raggedright

\bibitem{chani2016ilc}
J. Chani-Cahuana, P. N. Landin, C. Fager, and T. Eriksson,
``Iterative learning control for RF power amplifier linearization,''
\emph{IEEE Transactions on Microwave Theory and Techniques}, vol. 64, no. 9,
pp. 2778-2789, 2016. \href{https://doi.org/10.1109/TMTT.2016.2588483}{doi: 10.1109/TMTT.2016.2588483}.

\bibitem{schoukens2017preinverse}
M. Schoukens, J. Hammenecker, and A. Cooman,
``Obtaining the preinverse of a power amplifier using iterative learning control,''
\emph{IEEE Transactions on Microwave Theory and Techniques}, vol. 65, no. 11,
pp. 4266-4273, 2017. \href{https://doi.org/10.1109/TMTT.2017.2694822}{doi: 10.1109/TMTT.2017.2694822}.

\bibitem{morgan2006gmp}
D. R. Morgan, Z. Ma, J. Kim, M. G. Zierdt, and J. Pastalan,
``A generalized memory polynomial model for digital predistortion of RF power amplifiers,''
\emph{IEEE Transactions on Signal Processing}, vol. 54, no. 10,
pp. 3852-3860, 2006. \href{https://doi.org/10.1109/TSP.2006.879264}{doi: 10.1109/TSP.2006.879264}.

\bibitem{tarver2019backprop}
C. Tarver, L. Jiang, A. Sefidi, and J. R. Cavallaro,
``Neural network DPD via backpropagation through a neural network model of the PA,''
in \emph{Proc. 53rd Asilomar Conference on Signals, Systems, and Computers},
pp. 358-362, 2019. \href{https://doi.org/10.1109/IEEECONF44664.2019.9048910}{doi: 10.1109/IEEECONF44664.2019.9048910}.

\bibitem{loebl2023memorybackprop}
E. Loebl, N. Ginzberg, and E. Cohen,
``Direct learning neural network digital predistortion using backpropagation through a memory power amplifier model,''
in \emph{Proc. IEEE/MTT-S International Microwave Symposium},
pp. 791-794, 2023. \href{https://doi.org/10.1109/IMS37964.2023.10187912}{doi: 10.1109/IMS37964.2023.10187912}.

\bibitem{wei2026icg}
X. Wei, Y. Liu, W. Pan, W. Ma, Q. Xu, S. Shao, and Y. Tang,
``Iterative learning for RF power amplifier linearization based on instantaneous complex gain,''
\emph{IEEE Microwave and Wireless Technology Letters}, vol. 36, no. 1,
pp. 7-10, 2026. \href{https://doi.org/10.1109/LMWT.2025.3620316}{doi: 10.1109/LMWT.2025.3620316}.

\end{thebibliography}

\end{document}
